\documentclass[11pt]{article}

\usepackage[margin=1in]{geometry}
\usepackage{amsmath,amssymb,amsthm}
\usepackage{mathtools}
\usepackage{bm}
\usepackage{hyperref}
\usepackage{booktabs}
\usepackage{graphicx}
\usepackage{enumitem}
\usepackage{microtype}
\usepackage{courier}
\usepackage[T1]{fontenc}

\title{Multiplicity Theory, Recursive Becoming, and Prime-Locked EEG:\\
A Unified Computational Ontology with an Empirical Pipeline}

\author{Citizen Gardens \\ \\ The Foundation of Multiplicity \\
\vspace{0.3em}
\small April 3, 2026}

\date{}

\begin{document}

\maketitle

\begin{abstract}
This report consolidates recent developments on the ``13+1 Strata of Recursive Becoming''---a unified computational ontology that seeks to connect mathematics, physics, computation, consciousness, and social dynamics into a single recursive framework.\footnote{Original ontology draft: \emph{A Unified Computational Ontology: The 13+1 Strata of Recursive Becoming}.} At the core of this architecture lies a recursive tensor field \(\Xi(t)\) governed by a universal multiplicity constant \(\Lambda_m\), a multiplicative operator \(M\), and a torsion-like feedback term. Building on this conceptual foundation, we introduce a rigorous multiplicity functor \(\mathfrak{M}\) that maps stratum-specific states into prime-indexed multiplicity spectra, and we instantiate \(\mathfrak{M}_2\) concretely for EEG data (the ``qualia'' stratum). We provide a preregisterable MNE-Python pipeline that transforms raw EEG recordings into prime-locked multiplicity distributions using Welch spectral estimates, robust channel weighting, and regional pooling. The report includes a comprehensive mathematical overview, the definition of the multiplicity category, a minimal three-stratum architecture, and a code skeleton implementing \(\mathfrak{M}_2\) and the coherence statistic based on Jensen--Shannon distance. This document is intended to serve both as a conceptual backbone for Multiplicity Theory and as an executable plan for the first empirical tests of cross-stratum multiplicity coherence.
\end{abstract}
\newpage
\tableofcontents
\newpage
\section{Executive Summary}

The ``13+1 Strata of Recursive Becoming'' framework proposes a maximal unification of atomic, biological, cognitive, social, and cosmological dynamics within a single computational ontology.\footnote{See Sections 1--3 of the original ontology draft for the high-level description of the strata.} At its core is a recursive tensor field
\begin{equation}
  \frac{d\Xi(t)}{dt} \;=\; \Lambda_m M \Xi(t) + [M,\Xi(t)],
  \label{eq:recursive-tensor-dynamics}
\end{equation}
where \(\Lambda_m\) is a universal multiplicity constant, \(M\) a multiplicative operator, and \([M,\Xi(t)]\) a torsion term representing recursive feedback.

The main developments captured in this report are:

\begin{enumerate}[leftmargin=*]
  \item \textbf{Multiplicity Functor.} We formalize a functor \(\mathfrak{M} : \mathcal{C}_\Xi \to \mathcal{C}_{\mathrm{mult}}\) which sends stratum-specific states (adelic fields, qualia operads, social density matrices, etc.) to prime-indexed multiplicity spectra. The target category \(\mathcal{C}_{\mathrm{mult}}\) consists of probability distributions over a finite set of primes, with morphisms given by Markov operators that are Lipschitz under a chosen metric (e.g., Jensen--Shannon or Wasserstein-1).
  \item \textbf{Minimal Working Triad.} We identify a minimal three-stratum stack:
  \begin{itemize}
    \item Stratum 0: a finite adelic toy model over a small prime set \(P_0\),
    \item Stratum 2: a qualia operad realized as EEG time windows,
    \item Stratum 4: quantum social coherence realized as agent-based networks.
  \end{itemize}
  Each stratum is equipped with its own multiplicity functor \(\mathfrak{M}_0, \mathfrak{M}_2, \mathfrak{M}_4\), and cross-stratum coherence is defined as bounded distance in multiplicity space.
  \item \textbf{Concrete EEG Pipeline for \(\mathfrak{M}_2\).} We provide a preregisterable MNE-Python pipeline that maps raw 64-channel EEG data into prime-locked multiplicity spectra:
  \begin{itemize}
    \item Prime set \(P = \{2,3,5,7,11,13\}\), base frequency \(f_0 = 4 \text{ Hz}\),
    \item Half-octave frequency bins centered at \(f_p = f_0 \log p\),
    \item Welch PSD with 2 s segments and 50\% overlap in the 1--30 Hz band,
    \item Robust channel weighting via inverse variance and bad-channel exclusion,
    \item Regional pooling (frontal, central, parietal, occipital),
    \item Normalization to probability distributions over \(P\).
  \end{itemize}
  \item \textbf{Coherence Statistic.} We define cross-stratum coherence via Jensen--Shannon distance between multiplicity spectra from the qualia stratum and a parallel social stratum \(\mathfrak{M}_4\). A decrease in this distance under conditions of increased coupling is the primary empirical signature of multiplicity coherence.
\end{enumerate}

Together, these components turn the original manifesto into a falsifiable research program:

\begin{itemize}
  \item The ontology is anchored by invariants (multiplicity spectra) and explicit functors.
  \item The minimal triad and EEG pipeline are ready for numerical sanity checks and preregistration.
  \item A subsequent stability lemma will link Lipschitz dynamics in multiplicity space to persistent cross-stratum coherence.
\end{itemize}

\section{Mathematical Foundations}

\subsection{Recursive Tensor Field Dynamics}

The original framework introduces a recursive tensor field \(\Xi(t)\) as a multi-scalar object governing dynamics across scales:
\begin{equation}
  \frac{d\Xi(t)}{dt} \;=\; \Lambda_m M \Xi(t) + [M,\Xi(t)], \tag{\ref{eq:recursive-tensor-dynamics}}
\end{equation}
where:
\begin{itemize}
  \item \(\Lambda_m\) is the \emph{Universal Multiplicity Constant}, stabilizing recursive interactions.
  \item \(M\) is a \emph{Multiplicative Operator} that modulates recursive strain.
  \item \([M,\Xi(t)]\) is a \emph{recursive torsion} term, encoding nonlinear feedback loops.
\end{itemize}

In addition, a prime-indexed tensor field
\begin{equation}
  \Psi_p(t) \;=\; \sum_{i=1}^N p_i \, T(p_i,t)
\end{equation}
is introduced to bridge discrete prime structure and continuous dynamics. The report refines this idea by treating \(\Psi_p(t)\) not as a single scalar but as the \emph{shadow} of a richer multiplicity functor \(\mathfrak{M}\).

\subsection{Strata and Categories}

We abstract the multi-stratum ontology into a category \(\mathcal{C}_\Xi\):

\begin{itemize}
  \item Objects in \(\mathcal{C}_\Xi\) are stratum-specific state spaces such as:
  \begin{align*}
    \Xi_{\text{adelic}}(t) &\in \mathbf{Adel}_{P_0},\\
    \Xi_{\text{qualia}}(t) &\in \mathbf{Oper}_{\text{EEG}},\\
    \Xi_{\text{social}}(t) &\in \mathbf{ProbNet}_N,
  \end{align*}
  where \(\mathbf{Adel}_{P_0}\) is a finite adelic module over a prime set \(P_0\), \(\mathbf{Oper}_{\text{EEG}}\) a category of operadic structures on EEG windows, and \(\mathbf{ProbNet}_N\) a category of social networks with probability distributions on messages.
  \item Morphisms in \(\mathcal{C}_\Xi\) are recursive updates (e.g., time-evolution maps) and cross-stratum functors (e.g., projection from adelic states to EEG-relevant quantities).
\end{itemize}

The global dynamics can be viewed as an endofunctor
\[
  T : \mathcal{C}_\Xi \to \mathcal{C}_\Xi
\]
such that, at the level of objects, it reduces to equation~\eqref{eq:recursive-tensor-dynamics}.

\subsection{Multiplicity Category \(\mathcal{C}_{\mathrm{mult}}\)}

We fix a finite set of primes \(P = \{p_1, \dots, p_k\}\). Let \(\Delta(P)\) denote the simplex of probability distributions on \(P\), i.e.
\[
  \Delta(P) = \left\{ w : P \to [0,1] \,\bigg|\, \sum_{p \in P} w(p) = 1 \right\}.
\]

We define the \emph{multiplicity category} \(\mathcal{C}_{\mathrm{mult}}\) as follows:

\begin{itemize}
  \item Objects: elements \(w \in \Delta(P)\), interpreted as prime-indexed multiplicity spectra.
  \item Morphisms: Markov operators \(K : \Delta(P) \to \Delta(P)\) of the form
  \[
    (Kw)(p_j) = \sum_i K_{ji} w(p_i),
  \]
  where \(K_{ji} \geq 0\) and \(\sum_j K_{ji} = 1\) for each \(i\).
  \item Metric: a fixed metric \(d\) on \(\Delta(P)\), e.g.\ Jensen--Shannon distance or Wasserstein-1. Morphisms of interest are 1-Lipschitz with respect to \(d\).
\end{itemize}

In practice, Jensen--Shannon distance is especially convenient because it is symmetric, bounded, and provided directly in SciPy's distance module.\footnote{See \texttt{scipy.spatial.distance.jensenshannon} for the implementation.}

\subsection{Multiplicity Functor \(\mathfrak{M}\)}

We now define the multiplicity functor(s) \(\mathfrak{M}\).

\begin{definition}[Multiplicity Functor]
Let \(\mathcal{C}_\Xi\) be the category of stratum states and \(\mathcal{C}_{\mathrm{mult}}\) the multiplicity category defined above. A \emph{multiplicity functor} is a (covariant) functor
\[
  \mathfrak{M} : \mathcal{C}_\Xi \to \mathcal{C}_{\mathrm{mult}}
\]
that assigns to each state its prime-indexed multiplicity spectrum and to each morphism a Markov operator on \(\Delta(P)\).
\end{definition}

In practice, we define stratum-specific functors:
\[
  \mathfrak{M}_0 : \mathcal{A}_0 \to \mathcal{C}_{\mathrm{mult}},\quad
  \mathfrak{M}_2 : \mathcal{A}_2 \to \mathcal{C}_{\mathrm{mult}},\quad
  \mathfrak{M}_4 : \mathcal{A}_4 \to \mathcal{C}_{\mathrm{mult}},
\]
for the adelic, qualia, and social strata respectively, and then glue them into a single \(\mathfrak{M}\) on \(\mathcal{C}_\Xi\) by forming coproducts or products of categories.

\subsection{Multiplicities for the Minimal Triad}

\paragraph{Stratum 0: Adelic Toy Model.}
Take \(P_0 \subseteq P\) and consider states
\[
  \Xi_{\text{adelic}} \in \prod_{p \in P_0} \mathbb{Q}_p^{n_p}.
\]
Define a local norm \(\|\Xi_{\text{adelic}}\|_p\) at each prime, for example by the sum of p-adic norms of coordinates, and set
\[
  w_{\text{adelic}}(p) = \frac{\|\Xi_{\text{adelic}}\|_p}{\sum_{q\in P_0} \|\Xi_{\text{adelic}}\|_q},\quad p \in P_0.
\]
This yields a probability distribution on \(P_0\), zero-extended to \(P\).

\paragraph{Stratum 2: Qualia Operad (EEG).}
Here, \(\Xi_{\text{qualia}}(t)\) is realized as multi-channel EEG time windows. The construction of \(\mathfrak{M}_2\) is detailed in Section~\ref{sec:m2-eeg}, where PSD estimates are integrated over prime-indexed frequency bins to yield multiplicity spectra \(w_{\text{qualia}}\).

\paragraph{Stratum 4: Social Coherence.}
States \(\Xi_{\text{social}}(t)\) are small agent networks with message distributions. A natural multiplicity map uses eigenvalues of a covariance or adjacency-derived matrix: sorted eigenvalues \(\lambda_r\) are mapped to primes \(p_r\), and
\[
  w_{\text{social}}(p_r) = \frac{\lambda_r}{\sum_i \lambda_i}.
\]

\subsection{Cross-Stratum Coherence}

Let \(w_i(t) = \mathfrak{M}_i(\Xi_i(t))\) for \(i \in \{0,2,4\}\). Cross-stratum coherence is defined by inequalities of the form
\begin{align}
  d\big(w_0(t), w_2(t)\big) &\leq \varepsilon,\\
  d\big(w_2(t), w_4(t)\big) &\leq \varepsilon,
\end{align}
for some tolerance \(\varepsilon > 0\). The goal of the dynamics~\eqref{eq:recursive-tensor-dynamics} is to evolve \(\Xi(t)\) such that these inequalities hold over extended time intervals. Empirically, \(d\) will be instantiated as Jensen--Shannon distance between multiplicity spectra.

A forthcoming stability lemma will show that if the per-stratum update maps are Lipschitz in multiplicity space, then cross-stratum multiplicity distances remain bounded under the recursive dynamics.

\section{EEG-Based Multiplicity Functor \texorpdfstring{\(\mathfrak{M}_2\)}{M2}}
\label{sec:m2-eeg}

\subsection{Design Decisions}

We specify \(\mathfrak{M}_2\) as a fully determined pipeline from raw EEG data to multiplicity spectra.

\paragraph{Prime Set and Base Frequency.}
We fix the prime set
\[
  P = \{2,3,5,7,11,13\}
\]
and define prime-centered frequencies
\[
  f_p = f_0 \log p,\quad f_0 = 4 \text{ Hz}.
\]
This places all \(f_p\) approximately in the 2.5--11 Hz band, covering theta and alpha rhythms.

\paragraph{Frequency Bins.}
For each prime \(p\in P\), define a half-octave bin
\[
  \mathrm{bin}(p) = \left[ \frac{f_p}{\sqrt{2}}, f_p \sqrt{2} \right],
\]
and restrict analysis to the global frequency range 1--30 Hz.

\subsection{Preprocessing Pipeline}

We assume standard 64-channel EEG with 10--20 montage and sampling frequency \(\geq 500\) Hz.

\begin{enumerate}[label=\arabic*., leftmargin=*]
  \item \textbf{Data Loading and Montage.} Read raw EEG data (e.g., BIDS format) and apply a standard 10--20 montage.
  \item \textbf{Filtering.} Apply a 1--30 Hz bandpass filter (zero-phase FIR) and, if necessary, a notch filter at 50 or 60 Hz plus relevant harmonics.
  \item \textbf{Resampling.} Optionally resample to a common sampling frequency (e.g., 500 Hz) to standardize PSD parameters.
  \item \textbf{Rereferencing.} Apply average reference after marking bad channels.
  \item \textbf{Epoching.}
  \begin{itemize}
    \item Resting-state: non-overlapping 4-second epochs.
    \item Task/interaction: 4-second sliding windows with 50\% overlap within relevant blocks.
  \end{itemize}
  \item \textbf{Artifact Rejection.}
  \begin{itemize}
    \item Mark bad channels by visual inspection, high-variance criteria, or automated routines.
    \item Reject epochs where any channel exceeds 150~\(\mu\)V peak-to-peak amplitude.
    \item Optionally run ICA to remove ocular artifacts, with parameters fixed a priori.
  \end{itemize}
\end{enumerate}

\subsection{PSD Estimation}

For each subject and each clean epoch \(E\), we compute channel-wise PSDs using Welch's method via MNE-Python:
\begin{itemize}
  \item PSD function: \(\texttt{epochs.compute\_psd(method="welch")}\),
  \item Frequency range: \(f_{\mathrm{min}} = 1.0\) Hz, \(f_{\mathrm{max}} = 30.0\) Hz,
  \item Segment length: \(n_{\mathrm{per\_seg}} = 2.0 \times \texttt{sfreq}\) (2 s segments),
  \item Overlap: 50\% (i.e., \(n_{\mathrm{overlap}} = n_{\mathrm{per\_seg}}/2\)).
\end{itemize}
These choices align with standard spectral analysis practices in MNE-Python.\footnote{See MNE's documentation for \texttt{psd\_welch} and the tutorials on sensor-space spectral analysis.}

Denote by \(S_c(f;E)\) the PSD for channel \(c\) in epoch \(E\).

\subsection{Channel Weights and Robustness}

To downweight noisy channels while respecting bad-channel annotations, we define weights \(\omega_c\) as follows:
\begin{enumerate}[leftmargin=*]
  \item Compute channel-wise variance \(\sigma_c^2\) across all epochs and time samples.
  \item For channels marked as bad in \(\texttt{info['bads']}\), set \(\sigma_c^2 = +\infty\).
  \item Define inverse-variance weights
  \[
    \omega_c = \frac{\sigma_c^{-2}}{\sum_{d} \sigma_d^{-2}},
  \]
  treating infinite variances as zero weights.
\end{enumerate}
These weights are fixed per subject and used when aggregating PSD across channels.

\subsection{Global Multiplicity Spectrum}

For each epoch \(E\), we define the global multiplicity spectrum \(w_E\) by:

\begin{enumerate}[leftmargin=*]
  \item For each prime \(p \in P\), let \(A_p(E)\) be the weighted PSD mass in \(\mathrm{bin}(p)\):
  \[
    A_p(E) = \sum_{c} \omega_c \int_{\mathrm{bin}(p)} S_c(f;E) \, df.
  \]
  Numerically, this is computed by summing PSD values over the frequency indices falling into the bin.
  \item Normalize to obtain a probability distribution:
  \[
    w_E(p) = \frac{A_p(E)}{\sum_{q \in P} A_q(E)}.
  \]
\end{enumerate}

Thus, the global multiplicity functor \(\mathfrak{M}_2\) is realized as
\[
  \mathfrak{M}_2(E) = w_E \in \Delta(P).
\]

\subsection{Regional Multiplicity Spectra}

Let \(R\) be a set of anatomical regions (e.g., frontal, central, parietal, occipital) with fixed channel lists. For each region \(r\), define region-specific weights \(\omega_c^{(r)}\), e.g., uniform within region or inverse-variance restricted to the region.

Then for each region \(r\) and epoch \(E\), define
\[
  A_p^{(r)}(E) = \sum_{c \in r} \omega_c^{(r)} \int_{\mathrm{bin}(p)} S_c(f;E)\,df,
\]
and
\[
  w_E^{(r)}(p) = \frac{A_p^{(r)}(E)}{\sum_{q\in P} A_q^{(r)}(E)}.
\]

This yields a regional multiplicity spectrum \(\mathfrak{M}_2^{(r)}(E) = w_E^{(r)}\) for each region, enabling topographic analyses of multiplicity structure.

\subsection{Operadic Composition}

At the level of multiplicity spectra, the operadic operations of time concatenation and region pooling are implemented as convex combinations:

\paragraph{Time Concatenation.}
Let \(E_1, \dots, E_n\) be epochs of durations \(T_i\) merged into a composite window \(E^\ast\). Define
\[
  w_{E^\ast}(p) = \frac{\sum_{i} T_i w_{E_i}(p)}{\sum_i T_i}.
\]

\paragraph{Region Pooling.}
Let \(r_1, r_2\) be regions merged into \(r^\ast\) with mixing coefficient \(\alpha\). Set
\[
  w_{E}^{(r^\ast)}(p) = \alpha \, w_{E}^{(r_1)}(p) + (1-\alpha)\, w_{E}^{(r_2)}(p).
\]

These constructions are clearly Markov (stochastic linear maps) and thus define valid morphisms in \(\mathcal{C}_{\mathrm{mult}}\).

\section{Coherence with Social Stratum}

Let \(w_{\text{qualia}}(t) = \mathfrak{M}_2(\Xi_{\text{qualia}}(t))\) and \(w_{\text{social}}(t) = \mathfrak{M}_4(\Xi_{\text{social}}(t))\) be the multiplicity spectra for EEG and social data in matched time windows.

\subsection{Jensen--Shannon Coherence Statistic}

We define the global coherence at time \(t\) as the Jensen--Shannon distance
\[
  d_{\mathrm{JS}}\big(w_{\text{qualia}}(t), w_{\text{social}}(t)\big),
\]
and similarly regional coherence via
\[
  d_{\mathrm{JS}}\big(w_{\text{qualia}}^{(r)}(t), w_{\text{social}}^{(r)}(t)\big).
\]

The primary hypothesis is that under increased coupling (e.g., interactive social tasks versus rest), these distances decrease, indicating cross-stratum multiplicity alignment:
\[
  \mathbb{E}\big[d_{\mathrm{JS}}^{\text{rest}}\big] \;>\; \mathbb{E}\big[d_{\mathrm{JS}}^{\text{interaction}}\big].
\]

\subsection{Minimal \texorpdfstring{\(\mathfrak{M}_4\)}{M4} Placeholder}

A concrete \(\mathfrak{M}_4\) can be instantiated by:
\begin{itemize}
  \item Constructing a covariance matrix of message counts or probabilities across agents.
  \item Computing eigenvalues \(\lambda_r\), sorting them, and mapping rank \(r\) to prime \(p_r\).
  \item Defining
  \[
    w_{\text{social}}(p_r) = \frac{\lambda_r}{\sum_i \lambda_i}.
  \]
\end{itemize}
This parallel construction ensures that both EEG and social layers live in the same multiplicity space \(\Delta(P)\).

\section{Code Skeleton for \texorpdfstring{\(\mathfrak{M}_2\)}{M2} and Coherence}

This section presents a condensed version of the Python/MNE code skeleton implementing the EEG multiplicity pipeline. The code assumes MNE-Python, NumPy, and SciPy are installed.

\subsection{Constants and Prime Bins}

\begin{verbatim}
import mne
import numpy as np
from scipy.spatial.distance import jensenshannon

PRIMES = np.array([2, 3, 5, 7, 11, 13], dtype=float)
F0 = 4.0  # Hz

FMIN, FMAX = 1.0, 30.0
N_PER_SEG_SECONDS = 2.0

FP = F0 * np.log(PRIMES)
BIN_LOW = FP / np.sqrt(2.0)
BIN_HIGH = FP * np.sqrt(2.0)

REGIONS = {
    "frontal":   ["Fp1", "Fp2", "F3", "F4", "F7", "F8", "Fz"],
    "central":   ["C3", "C4", "Cz"],
    "parietal":  ["P3", "P4", "Pz"],
    "occipital": ["O1", "O2", "Oz"],
}
\end{verbatim}

\subsection{PSD and Channel Weights}

\begin{verbatim}
def compute_epoch_psd(epochs):
    sfreq = epochs.info["sfreq"]
    n_per_seg = int(N_PER_SEG_SECONDS * sfreq)
    psd = epochs.compute_psd(
        method="welch",
        fmin=FMIN,
        fmax=FMAX,
        n_per_seg=n_per_seg,
        n_overlap=n_per_seg // 2,
        verbose=False,
    )
    freqs = psd.freqs
    psd_data = psd.get_data()  # (n_epochs, n_channels, n_freqs)
    return freqs, psd_data

def compute_channel_weights(epochs):
    data = epochs.get_data()  # (n_epochs, n_channels, n_times)
    var = data.reshape(-1, data.shape[2]).var(axis=1)
    var = var.reshape(data.shape[0], data.shape[1]).mean(axis=0)
    bads = set(epochs.info.get("bads", []))
    ch_names = epochs.ch_names
    for idx, name in enumerate(ch_names):
        if name in bads:
            var[idx] = np.inf
    inv_var = 1.0 / (var + 1e-12)
    if np.isinf(inv_var).all():
        inv_var = np.ones_like(inv_var)
    weights = inv_var / inv_var.sum()
    return weights  # (n_channels,)
\end{verbatim}

\subsection{Global and Regional Multiplicity Spectra}

\begin{verbatim}
def compute_w_qualia(freqs, psd_data, weights):
    n_epochs, _, _ = psd_data.shape
    n_primes = len(PRIMES)
    w_global = np.zeros((n_epochs, n_primes))
    idx_bins = [
        np.where((freqs >= low) & (freqs <= high))[0]
        for low, high in zip(BIN_LOW, BIN_HIGH)
    ]
    for e in range(n_epochs):
        epoch_psd = psd_data[e]
        A_p = np.zeros(n_primes)
        for j, idx in enumerate(idx_bins):
            if idx.size == 0:
                continue
            band_power = epoch_psd[:, idx].sum(axis=1)
            A_p[j] = np.dot(weights, band_power)
        denom = A_p.sum()
        w_global[e] = A_p / denom if denom > 0 else np.ones(n_primes)/n_primes
    return w_global

def compute_w_qualia_regional(freqs, psd_data, epochs):
    ch_names = epochs.ch_names
    n_epochs = psd_data.shape[0]
    n_primes = len(PRIMES)
    idx_bins = [
        np.where((freqs >= low) & (freqs <= high))[0]
        for low, high in zip(BIN_LOW, BIN_HIGH)
    ]
    region_spectra = {}
    for region, chans in REGIONS.items():
        idx_ch = [ch_names.index(ch) for ch in chans if ch in ch_names]
        if not idx_ch:
            continue
        weights_region = np.ones(len(idx_ch))
        weights_region = weights_region / weights_region.sum()
        w_reg = np.zeros((n_epochs, n_primes))
        for e in range(n_epochs):
            epoch_psd = psd_data[e, idx_ch, :]
            A_p = np.zeros(n_primes)
            for j, idx in enumerate(idx_bins):
                if idx.size == 0:
                    continue
                band_power = epoch_psd[:, idx].sum(axis=1)
                A_p[j] = np.dot(weights_region, band_power)
            denom = A_p.sum()
            w_reg[e] = A_p / denom if denom > 0 else np.ones(n_primes)/n_primes
        region_spectra[region] = w_reg
    return region_spectra
\end{verbatim}

\subsection{Coherence via Jensen--Shannon Distance}

\begin{verbatim}
def js_distance_matrix(w_a, w_b):
    assert w_a.shape == w_b.shape
    n_epochs = w_a.shape[0]
    d_js = np.zeros(n_epochs)
    for e in range(n_epochs):
        d_js[e] = jensenshannon(w_a[e], w_b[e])
    return d_js
\end{verbatim}

\subsection{End-to-End Subject Pipeline}

\begin{verbatim}
def run_subject_pipeline(raw_path, events=None, event_id=None,
                         tmin=0.0, tmax=4.0):
    raw = mne.io.read_raw(raw_path, preload=True)
    raw.set_montage("standard_1020")
    if raw.info["sfreq"] != 500.0:
        raw.resample(500.0)
    raw.filter(1.0, 30.0, fir_design="firwin")
    raw.notch_filter(freqs=[50.0])

    if events is None:
        events = mne.make_fixed_length_events(raw, id=1, duration=4.0)
        event_id = dict(rest=1) if event_id is None else event_id

    epochs = mne.Epochs(
        raw, events=events, event_id=event_id,
        tmin=tmin, tmax=tmax, baseline=None, preload=True
    )
    reject_criteria = dict(eeg=150e-6)
    epochs.drop_bad(reject=reject_criteria)

    freqs, psd_data = compute_epoch_psd(epochs)
    weights = compute_channel_weights(epochs)
    w_global = compute_w_qualia(freqs, psd_data, weights)
    w_regions = compute_w_qualia_regional(freqs, psd_data, epochs)
    return epochs, w_global, w_regions
\end{verbatim}

This skeleton can be extended with a concrete \(\mathfrak{M}_4\) and social data ingestion to compute full cross-stratum coherence.

\section{Outlook}

The developments summarized here mark a transition from a high-level unification manifesto to a concrete, falsifiable research program. The next steps are:

\begin{itemize}
  \item Prove a stability lemma showing that Lipschitz dynamical updates in multiplicity space preserve bounded cross-stratum distances.
  \item Instantiate \(\mathfrak{M}_4\) concretely on real social interaction data (e.g., chat logs, behavioral events) and integrate it with the EEG pipeline.
  \item Run pilot analyses on public EEG datasets (e.g., resting-state or interaction datasets from OpenNeuro) to test the reliability and discriminability of prime-locked multiplicity spectra.
\end{itemize}

If successful, this line of work will supply both conceptual and empirical evidence that prime-indexed multiplicity is a useful invariant for describing recursive becoming across physical, neural, and social strata.

\nocite{*}
\bibliographystyle{plainnat}
\bibliography{references} 
\end{document}